% =====================================================================
\section{Capstone: Cross-Sectional Test of U.S.\ Sector ETFs}
% =====================================================================

\subsection{Step 1: Dataset}

Of the datasets offered, this capstone uses \textbf{U.S.\ sector ETF returns}.

\subsection{Research Question}

Does market beta explain cross-sectional differences in average returns across
U.S.\ sector ETFs?

\subsection{Hypothesis}

The dependent variable is the average monthly return of each U.S.\ sector ETF,
and the main independent variable is its market beta relative to SPY.

I expect sector ETFs with higher market betas to have higher average returns.
Under the CAPM, the relationship should be positive, with the coefficient on
beta approximately reflecting the market risk premium. Economically, investors
should require higher expected returns for bearing greater systematic market
risk.

I would view the results as supportive of the CAPM if the beta coefficient is
positive and statistically significant, with an intercept reasonably close to
the risk-free rate. If the beta coefficient is insignificant, zero, or
negative, or if the intercept differs substantially from the risk-free rate, I
would conclude that the CAPM does not explain the cross-section of sector ETF
returns well.

\paragraph{Specification clarification (added after analysis).} The implemented
regression uses excess returns rather than raw returns. Therefore, the
Sharpe-Lintner CAPM predicts an intercept of zero, rather than the risk-free
rate. The original hypothesis above is retained as written before examining the
data.

\subsection{Methodology}

I use monthly returns for nine U.S.\ sector ETFs (XLB, XLE, XLF, XLI, XLK, XLP,
XLU, XLV, and XLY), with SPY as the market proxy and the risk-free rate from
the Kenneth French Data Library.

The sample covers February 2000 through December 2025, with 311 monthly
observations.

I first calculate excess returns by subtracting the monthly risk-free rate from
each ETF's return and SPY's return. I then estimate each ETF's market beta
using a time-series regression of ETF excess returns on SPY excess returns.

Finally, I run a cross-sectional OLS regression of average ETF excess returns
on their estimated market betas. This regression provides the coefficient
estimates. I use Fama-MacBeth (FMB) standard errors as the main basis for
inference because OLS standard errors treat the nine ETF residuals as
independent, although sector returns share common monthly shocks.

Each month, I regress the nine ETF excess returns on the same fixed full-sample
betas, then take the time-series average of the 311 monthly intercept and slope
estimates. The FMB standard error is the time-series standard deviation of the
monthly estimates divided by $\sqrt{T}$. Because the betas are the same every
month, the FMB coefficient averages equal the OLS coefficients exactly; only
the standard errors differ.

\subsection{Results}

\begin{table}[H]
\centering
\caption{Cross-sectional CAPM test on nine sector ETFs.}
\small
\begin{tabular}{lrrrrrrr}
\toprule
Variable & Coefficient & OLS & OLS & OLS & FMB & FMB & FMB \\
 & & Std.\ Error & $t$-stat & $p$-value & Std.\ Error & $t$-stat & $p$-value \\
\midrule
Intercept & 0.0045 & 0.0009 & 4.8158 & 0.0019 & 0.0028 & 1.6032 & 0.1099 \\
Beta      & 0.0020 & 0.0010 & 2.1220 & 0.0715 & 0.0037 & 0.5506 & 0.5823 \\
\bottomrule
\end{tabular}
\label{tab:cap}
\end{table}

\begin{tight}
\item Number of sector ETFs: 9
\item Sample period: February 2000--December 2025
\item Monthly return observations ($T$): 311
\item Cross-sectional $R^2$: 0.3915
\item Average SPY excess return: 0.0060
\item OLS $p$-values use $N - 2 = 7$ degrees of freedom; FMB $p$-values use
$T - 1 = 310$ degrees of freedom.
\end{tight}

All coefficients and standard errors are reported in decimal units per month.
For comparison with Part I, which reports returns in percent per month, the
estimates should be multiplied by 100.

\begin{table}[H]
\centering
\caption{Beta premium vs.\ market premium.}
\small
\begin{tabular}{lr}
\toprule
& Value \\
\midrule
Estimated beta premium ($\gamma_M$) & 0.0020 \\
Average SPY excess return & 0.0060 \\
Difference ($\gamma_M -$ SPY) & $-0.0040$ \\
Std.\ error of difference (monthly series) & 0.0027 \\
$t$-stat of difference & $-1.47$ \\
$p$-value of difference & 0.1416 \\
\bottomrule
\end{tabular}
\label{tab:capdiff}
\end{table}

The test of the difference uses the monthly series
$\gamma_{Mt} - \text{SPY}_t$, whose time-series standard error accounts for the
sampling uncertainty of both averages and their covariance. There is no
equivalent monthly series for the single cross-sectional OLS regression, so
under OLS the comparison is reported descriptively only.

\subsection{Step 4(a): Hypothesis}

The estimated beta premium is positive at $0.20\%$ per month, consistent with
my hypothesis, but insignificant using Fama-MacBeth standard errors
($p = 0.5823$). It is below the average SPY excess return of $0.60\%$, while
the estimated intercept is $0.45\%$ rather than zero, suggesting a flatter SML
than the CAPM predicts. Although OLS finds a significant intercept
($p = 0.0019$), the FMB intercept test ($p = 0.1099$) and the paired
market-premium gap test ($p = 0.1416$) are insignificant. By my Step 2
criterion, the hypothesis is not supported because the beta premium is
insignificant. However, an insignificant slope is not by itself a rejection of
the CAPM, and the FMB tests do not reject its intercept or slope restrictions.

\subsection{Step 4(b): Comparison with Parts I and II}

My estimated beta premium of $0.20\%$ per month compares with $0.035\%$ for the
49 industry portfolios and $-0.576\%$ for the 25 size/BE-ME portfolios. All
three are insignificant using FMB standard errors and below their respective
market premiums. However, Parts I and II find significantly positive intercepts
and significant differences between their estimated beta premiums and market
premiums, while my ETF test does not. Part II also finds a significant
book-to-market premium after controlling for beta. The slope is positive for
Part I and my Capstone and negative for Part II. So the ETF cross-section looks
similar to US equities (a flat SML), but my test gives much less precise
evidence against the CAPM.

\subsection{Step 4(c): Main Limitation}

The main limitation is the small cross-section of nine sector ETFs. Their
relatively similar market exposures limit variation in beta, reducing the
test's statistical power. FMB standard errors are substantially larger than OLS
standard errors, and the estimated beta premium cannot be statistically
distinguished from the market premium. In addition, using estimated rather than
true betas can bias the slope toward zero, while SPY is only a proxy for the
true market portfolio.

\subsection{Step 3: AI Workflow Log}

\subsubsection*{Step 3(a): Data Loading and Cleaning}

\paragraph{Prompt (Close Paraphrase).} I asked the LLM to generate Python code
to download monthly adjusted-price data for nine U.S.\ sector ETFs (XLB, XLE,
XLF, XLI, XLK, XLP, XLU, XLV, XLY) and SPY from January 2000 through December
2025, calculate monthly returns, align the series, and check for missing
observations.

\paragraph{LLM Output and My Correction.} The code successfully downloaded the
data and produced 311 monthly return observations from February 2000 through
December 2025.

However, it checked for missing values only after dropping missing
observations, making the diagnostic uninformative. I revised the code to report
missing values before and after cleaning and explicitly used
\texttt{fill\_method=None} when calculating returns. The final dataset
contained no missing prices, so no observations were removed.

\subsubsection*{Step 3(b): Cross-Sectional Regression}

\subparagraph*{Initial OLS Regression}

\paragraph{Prompt (Close Paraphrase).} Please generate Python code to estimate
the market beta of each of the nine U.S.\ sector ETFs using SPY as the market
proxy. Then run a cross-sectional OLS regression of each ETF's average monthly
return on its estimated beta. Report the coefficients, standard errors,
$t$-statistics, $p$-values, and R-squared.

\paragraph{LLM Output and My Correction.} The LLM generated code to estimate
full-sample betas and regress average ETF returns on those betas. The initial
results were an intercept of $0.0061$, a beta coefficient of $0.0020$
($p = 0.0721$), and $R^2 = 0.3903$.

However, the code used raw rather than excess returns. I revised it to subtract
the monthly risk-free rate from both ETF and SPY returns, using data from the
Kenneth French Data Library. This matches the Sharpe-Lintner CAPM
specification, which predicts a zero intercept for excess returns.

\subparagraph*{Fama-MacBeth Robustness Check}

A later AI code review recommended Fama-MacBeth (FMB) standard errors because
the original cross-sectional OLS inference did not adequately account for
common shocks across sector returns.

\paragraph{Prompt (Close Paraphrase).} Keep my original cross-sectional OLS
regression as the main methodology. Add Fama-MacBeth standard errors only as a
robustness check, using the same fixed full-sample betas and monthly ETF excess
returns. Report the OLS and FMB results side by side, including coefficients,
standard errors, $t$-statistics and $p$-values, and compare the estimated beta
premium with the average SPY excess return. Do not add Shanken, Newey-West or
other extensions.

\paragraph{LLM Output and Verification.} The LLM generated monthly
cross-sectional regressions using fixed betas, then calculated FMB standard
errors from the time-series variation in the 311 monthly coefficient estimates.

I reran the script and verified that FMB and OLS produced identical coefficient
estimates but different standard errors, $t$-statistics, and $p$-values. With
FMB standard errors, neither the intercept ($p = 0.1099$) nor the beta
coefficient ($p = 0.5823$) is significant. I therefore use the FMB standard
errors as the main basis for inference.

\paragraph{Additional Correction: Market-Premium Gap Test.} The initial
comparison divided the gap by the beta premium's standard error alone, ignoring
uncertainty in the SPY average and its covariance with the estimated premium.

I replaced it with a test based on the monthly difference series between the
estimated beta premium and SPY excess returns. The gap is $-0.0040$, with
SE $= 0.0027$, $t = -1.47$ and $p = 0.1416$. The difference is not
statistically significant.

\subsubsection*{Step 3(c): Interpretation Assistance}

\subparagraph*{Initial Raw-Return Results}

\paragraph{Prompt.} I ran a cross-sectional CAPM test on nine U.S.\ sector ETFs
using monthly data from February 2000 to December 2025. The cross-sectional
regression of average monthly returns on estimated market beta produced an
intercept of $0.0061$ ($t = 6.47$, $p = 0.0003$), a beta coefficient of
$0.0020$ ($t = 2.12$, $p = 0.0721$), and an R-squared of $0.3903$. Please
interpret these results. Do they support my hypothesis that higher-beta sector
ETFs should have higher average returns, and how strong is the evidence for the
CAPM?

\paragraph{Summary of the LLM's Answer.} The LLM correctly interpreted the
positive beta coefficient as consistent with my hypothesis. It noted that the
coefficient was significant at the $10\%$ level but not at $5\%$, and that beta
explained approximately $39\%$ of the observed cross-sectional variation. It
concluded that the evidence for the CAPM was limited.

\paragraph{My Evaluation.} The LLM correctly interpreted the coefficient's
sign, magnitude, and significance. However, it overlooked the intercept and
failed to compare the estimated slope with the market risk premium, both of
which are central CAPM restrictions.

What it got wrong: it relied on OLS $p$-values without accounting for common
shocks across sector returns and did not emphasize the small cross-section or
estimated-beta uncertainty. Some omissions reflected my original prompt, which
did not provide the risk-free rate or market premium.

\subparagraph*{Re-assessment After Code Corrections}

Switching to excess returns reduced the intercept from $0.0061$ to $0.0045$,
while the beta coefficient remained approximately $0.0020$. The subsequent FMB
check showed that neither coefficient was significant, and the beta premium was
not significantly different from the SPY premium ($p = 0.1416$).

These additional checks came from a later AI code review, not from the original
LLM interpretation. They changed my initial interpretation of the significant
OLS intercept as evidence against the CAPM.

\subparagraph*{Final Excess-Return Results}

\paragraph{Prompt.} I ran a cross-sectional CAPM test using nine U.S.\ sector
ETFs, with SPY as the market proxy. My sample contains 311 monthly observations
from February 2000 through December 2025.

I estimated each ETF's beta by regressing its excess returns on SPY excess
returns. I then regressed average ETF excess returns on estimated beta. I also
calculated Fama-MacBeth standard errors using monthly cross-sectional
regressions with fixed betas.

Here are my final results (all coefficients and standard errors are in decimal
units per month):

\begin{center}
\small
\begin{tabular}{lrr}
\toprule
& Intercept & Beta \\
\midrule
Coefficient & 0.0045 & 0.0020 \\
OLS SE & 0.0009 & 0.0010 \\
OLS $t$-stat & 4.8158 & 2.1220 \\
OLS $p$-value & 0.0019 & 0.0715 \\
FMB SE & 0.0028 & 0.0037 \\
FMB $t$-stat & 1.6032 & 0.5506 \\
FMB $p$-value & 0.1099 & 0.5823 \\
\bottomrule
\end{tabular}
\end{center}

\noindent Cross-sectional R-squared: 0.3915\\
Number of ETFs: 9\\
Average monthly SPY excess return: 0.0060

The estimated beta premium minus the average SPY excess return is $-0.0040$. A
test based on the monthly difference series gives $t = -1.47$ and
$p = 0.1416$.

Please interpret these results. Do they support my hypothesis that higher-beta
sector ETFs earn higher average excess returns? What do they imply about the
CAPM, and how do the conclusions differ when using OLS versus Fama-MacBeth
standard errors?

\paragraph{Summary of the LLM's Answer.} The LLM correctly noted that the
positive beta coefficient is consistent with my hypothesis but insignificant at
the $5\%$ level under both OLS and FMB inference.

It explained that the positive intercept and relatively low beta premium
suggest a flatter SML than the CAPM predicts. Although the OLS intercept is
significant, neither coefficient is significant using FMB standard errors. The
gap between the beta premium and the SPY premium is also insignificant.

It concluded that the point estimates suggest a flatter SML but do not provide
sufficient evidence to reject the individually tested CAPM restrictions. It
also noted the small cross-section, estimated-beta uncertainty, potential
serial correlation, and absence of a formal joint test.

\paragraph{My Evaluation.} The LLM correctly distinguished OLS from FMB
inference and compared the estimated beta premium with the market premium
rather than simply testing whether beta is significant.

However, it did not fully explain the implications of the test's low precision.
The approximate $95\%$ FMB confidence interval for the beta premium is
$-0.53\%$ to $0.93\%$ per month, including both zero and the SPY premium of
$0.60\%$.

It also mentioned estimated-beta uncertainty without explaining that
measurement error can flatten the estimated SML. Finally, it overlooked Roll's
critique: SPY is only a proxy for the true market portfolio. I found no major
numerical error in its interpretation, but these limitations deserved more
attention.

\subsection{Step 5: AI Workflow Reflection}

Using an LLM helped me move quickly from my research design to working Python
code, allowing me to spend more time checking the results. I corrected the
missing-value diagnostics and changed the regression from raw to excess returns
to match the Sharpe-Lintner CAPM. A later AI review prompted me to add
Fama-MacBeth standard errors, which changed my interpretation of the
statistical evidence. These experiences showed me that AI can help both
generate code and question results, but I still need to verify its methods and
conclusions independently. Next time, I would specify the CAPM form, excess
returns, data checks, and inference method in my initial prompts.
